\begin{textsample}{POS Dim 4 – global\_south\_2023\_12 – Score 18.00 – global\_south\_2023\_12/104528226.txt}  \label{ex:f4_pos_147}
The UN security council caves to US pressure to prevent a \textbf{ceasefire} Israel and Hamas announced a \textbf{deal} on Wednesday allowing at least 50 hostages and scores of Palestinian \textbf{prisoners} to be \textbf{freed}, while offering besieged Gaza residents a \textbf{four-day} \textbf{truce} after weeks of all-out war.
The bottom line of the just-passed resolution at the UN Security Council is that it is not a \textbf{ceasefire} resolution.
It is not even a"suspension of hostilities"resolution, which reflected the first major concession to Washington ' s demands.
That would have turned the resolution into a repeat of last month ' s \textbf{temporary} pause -- potentially useful for allowing in some additional humanitarian aid, perhaps another **37 ; 1302 ; TOOLONG \textbf{swap}, and a few days \textbf{respite} for the millions of people in Gaza dying under Israeli bombardment before Israel ' s full-scale war started again.
But this resolution does not even do that.
Despite misleading headlines in way too much of the mainstream media, the only mention even of"humanitarian pauses"appears in reference to the council ' s November resolution that did call for such \textbf{temporary} halts to @ @ @ @ @ @ @ @ @ @ not anywhere in the operative paragraphs of the new resolution.
The operative paragraphs do not call for pausing, suspending, ending, easing or ceasing hostilities -- meaning Israel can continue its deadly assaults by air and land without violating the security council ' s fought-over resolution.
The vote was 13 in favour, with the US and Russia abstaining.
Moscow had \textbf{proposed} an amendment \textbf{returning} to the"suspension of hostilities"language, but despite 10 votes in favour and four abstentions, the amendment was \textbf{rejected} by a US veto.
Instead, the final text"calls for urgent steps to immediately allow safe, unhindered, and expanded humanitarian access"without defining those steps, and without any acknowledgement that the crucial"step"would require Israel to stop its bombing campaign and end its ground assaults.
That means Israel, the overwhelmingly stronger party responsible for the deaths of 20 000+ Palestinians, overwhelmingly children and women, can decide when or if its bombs, drones, tank assaults decimating the Gaza Strip and its @ @ @ @ @ @ @ @ @ @ So instead of calling for an actual \textbf{ceasefire}, or even a"suspension of hostilities,"the resolution calls for unnamed"urgent steps... to create the conditions for a sustainable \textbf{cessation} of hostilities".
This means that until Tel Aviv decides what steps it might want to take if it wants to create conditions for a \textbf{ceasefire} at all, the resolution means that the security council embraces the US position of allowing Israel ' s genocidal assault to continue unhindered.
The resolution demands that all parties"facilitate and enable the immediate, safe and unhindered delivery of humanitarian assistance at scale"directly to Palestinian civilians and that they"facilitate the use of all available routes to and throughout the entire Gaza strip"to deliver the desperately-needed aid.
If this were not happening at a moment in which the UN World Food Program is reporting that 90 \% of Gaza ' s two million + people are hungry and that half the population is starving, this would be laughable -- since it is clearly impossible to provide immediate and @ @ @ @ @ @ @ @ @ @ under Israel ' s unceasing bombs.
Another section demands that sufficient fuel be allowed into Gaza -- a good move in theory, since Israel had prohibited almost all fuel deliveries, but not worth much in the context of allowing continuing bombing raids across the strip.
Other sections request that the UN Secretary-General appoint someone to"oversee"the provision of aid -- but leaves Israel completely in control of the lethally slow inspection process that has kept hundreds of truckloads of water, food, medicine for the starving population stalled at the \textbf{Egyptian} side of the border.
And Israel will remain in control of the checkpoints and conditions on the ground within Gaza.
Earlier \textbf{drafts} called for a UN inspection regime to replace that of Israel.
But that language was stripped out.
The resolution demands the \textbf{release} of hostages in Gaza, which is good, but abandons any concern for the thousands of Palestinian \textbf{prisoners} held illegally in Israeli military \textbf{prisons} who would likely be \textbf{freed} in any \textbf{prisoner} \textbf{swap}.
Of course, since the actual @ @ @ @ @ @ @ @ @ @ underway outside the security council, that demand doesn't mean much.
Overall, the security council caved to US power.
Led by the UAE, the council ' s only Arab member, ostensibly representing other Arab countries in \textbf{negotiations} with Washington, most members of the council were excluded from the actual discussions leading to the weakened resolution.
The council could have better stuck to the principles of the early \textbf{draft} resolution, recognising the desperate need for a \textbf{ceasefire} -- and force the US to publicly acknowledge its isolation in the world by using its veto again.
That would have sent the issue back to the general assembly under special UN arrangements that allow a much greater level of enforcement than is generally true for general assembly resolutions.
But the concern with not antagonising the US ( the UAE remains a key US ally, as do Egypt, Qatar, Jordan, Saudi Arabia and others, after all ) was great enough that the council was willing to vote for an"aid resolution"that will do @ @ @ @ @ @ @ @ @ @, elders, women and men now being killed by US-made, Israeli-dropped bombs and by US-armed Israeli troops.
This is just so the US would not be embarrassed by having to use its veto again.
The repeated use of vetoes may cost the US government something at some point -- whether in domestic support, international legitimacy, or maybe even legally if Washington were to be held accountable for enabling Israel ' s genocidal assault.
But for now, if we ask who won this hard-fought battle in the security council, the answer is clear.
Not the Palestinian civilians whose lives were supposed to be protected by this resolution -- but the United States who now won't have to be embarrassed.
Phyllis Bennis is a fellow of the Institute for Policy Studies in Washington DC, and serves as the international adviser for Jewish Voice for Peace.

% matched lemmas: ceasefire, cessation, deal, draft, egyptian, four-day, free, negotiation, prison, prisoner, propose, reject, release, respite, return, swap, temporary, truce
\end{textsample}
